We review learning-based malware detection (Section~\ref{sec:XML:BgDetection}), post-hoc explanation
(Section~\ref{sec:XML:BgExplanation}) and the Shapley value it often rests on (Section~\ref{sec:XML:BgShapley}), how both meet in explaining malware detectors (Section~\ref{sec:XML:BgMalwareXAI}),
and prior ways of moving attributions between representations (Section~\ref{sec:XML:BgTransfer}), and close with where
explanation provenance fits (Section~\ref{sec:XML:BgPosition}).

\subsection{Learning-Based Malware Detection}
\label{sec:XML:BgDetection}

Static detectors classify a file from its contents, without running it. They differ above all in how they represent the
file. \emph{Feature-based} detectors compute hand-engineered features with a parser: byte $n$-grams in early
work~\cite{schultz2001data, kolter2004learning}, permissions and API calls for Android~\cite{arp2014drebin}, and, for
Windows (PE) executables, the EMBER feature set~\cite{anderson2018ember}, the standard for benchmarks such as
SOREL-20M~\cite{harang2020sorel} and BODMAS~\cite{yang2021bodmas}. Its third version~\cite{joyce2025ember} has 2{,}568
features in 12 groups: byte and byte-entropy histograms, string statistics, header fields, sections, imports and
exports hashed into buckets, data directories, the Rich header, the signature, parser warnings, and general file
properties. Its extractor is a program, not a model: it parses the file, reads fields, counts and hashes names, and
computes entropy over sliding windows. Gradient-boosted decision trees (GBDTs) such as LightGBM~\cite{ke2017lightgbm}
are the standard and strongest model family for such features. \emph{Raw-byte} detectors learn from the bytes
directly~\cite{raff2018malware, krvcal2018deep, raff2021classifying}: MalConv~\cite{raff2018malware} embeds each byte,
convolves over windows of 512 bytes, and pools over the whole file.

Deployed engines layer both kinds: Microsoft Defender combines models over featurized metadata with deep networks over
full file contents~\cite{microsoft2019insideout}, and the winning defense of a public evasion competition combined
detectors over distinct feature sets so that an attacker must evade all of them at once~\cite{quiring2020against}. The
separation between the file and its features also shapes attacks: a feature-space evasion must be realized as a valid
file, the inverse feature-mapping problem~\cite{pierazzi2020intriguing}, which practical attacks solve with
manipulations at known locations in the PE format~\cite{demetrio2021adversarial}. Learned detectors also pick up
spurious correlations from their training data~\cite{arp2022and}, which is one reason to explain them.

\subsection{Post-Hoc Explanation}
\label{sec:XML:BgExplanation}

An \emph{explainer} assigns each input of a model a share of its output, its \emph{attribution}. Perturbation methods
query the model on modified inputs: LIME~\cite{ribeiro2016should} fits a local surrogate, and
occlusion~\cite{zeiler2014visualizing} replaces one window of the input at a time. Shapley-value
methods~\cite{vstrumbelj2014explaining, lundberg2017unified} credit each input with its average marginal contribution
over coalitions of inputs, the absent inputs drawn from a background distribution; KernelSHAP estimates these values
for any model, and TreeSHAP~\cite{lundberg2020local} computes them exactly for tree ensembles, either
\emph{path-dependent} (absent features follow the training data down each tree) or \emph{interventional} (absent
features take values from a background set). Gradient methods such as Integrated
Gradients~\cite{sundararajan2017axiomatic} integrate a network's gradient along a path from a baseline to the input,
and propagation methods such as LRP~\cite{bach2015pixel} and DeepLIFT~\cite{shrikumar2017learning} redistribute the
output backward layer by layer. All of them attribute the inputs the model reads. Faithfulness is tested by deleting
the top-ranked inputs and measuring the drop in output~\cite{petsiuk2018rise}, by retraining without
them~\cite{hooker2019benchmark}, by sanity checks~\cite{adebayo2018sanity}, or against a planted
trigger~\cite{lin2021what}; explanations can also be manipulated~\cite{slack2020fooling}, a threat we do not address.

\subsection{Shapley Values}
\label{sec:XML:BgShapley}

Several of these explainers, and our construction of the provenance matrix (Section~\ref{sec:XML:Constructing}), rest
on the Shapley value from cooperative game theory~\cite{shapley1953value}. A \emph{game} over a set of players $N$ is a
function $v$ that assigns every coalition $S \subseteq N$ a value $v(S)$; $v(\emptyset)$ is the value of no players at
all. The Shapley value splits the total gain $v(N) - v(\emptyset)$ among the players. Its answer is each player's
average marginal contribution: add the players one at a time in a random order, record how much the value changes when
player $i$ joins, and average over all orders,
\[
\phi_i(v) = \sum_{S \subseteq N \setminus \{i\}} \frac{|S|! \, (|N| - |S| - 1)!}{|N|!} \bigl( v(S \cup \{i\}) - v(S) \bigr).
\]
It is the only split with four properties~\cite{shapley1953value}: the players' values sum to the total gain
(\emph{efficiency}); two players who add the same to every coalition receive the same value (\emph{symmetry}); a player
who never changes the value receives zero (\emph{null player}); and the values of a sum of two games are the sums of
their values (\emph{linearity}). The \emph{unanimity game} $u_T$ of a coalition $T$ shows them at work: $u_T(S) = 1$ if
$S$ contains all of $T$ and $0$ otherwise, so by symmetry and the null player each member of $T$ receives $1/|T|$ and
every other player $0$.

To explain a model, the players are its inputs, and $v(S)$ is the model's output when only the inputs in $S$ are
present. Explainers differ in what an absent input becomes: a fixed baseline, or values drawn from a background set, as
in interventional TreeSHAP and KernelSHAP. The Shapley values therefore explain the change in output from that baseline
to the input, and they change with it. Computing them exactly takes exponentially many coalitions in general, so
explainers sample coalitions or, like TreeSHAP, exploit the model's structure.

\subsection{Explaining Malware Detectors}
\label{sec:XML:BgMalwareXAI}

Security has its own requirements for explanations~\cite{guo2018lemna, warnecke2020evaluating, nadeem2023sok}, and
practitioners value them for verifying models and understanding security events, while finding existing methods
ill-matched to their tasks~\cite{mink2023everybody}, and the analyst's object, the file, and the detector's input, the
feature vector, rarely coincide~\cite{rui2024position}. Explanations follow the representation. Feature-based detectors
are explained at the feature level~\cite{melis2018explaining, fan2020can, wu2021android, zhu2024dynamic}, most often
with SHAP over EMBER features, whose reliability on its own has been questioned~\cite{mohseni2026shap}. Raw-input
detectors are explained over the input: activations and gradients show that MalConv relies heavily on the
header~\cite{coull2019activation, bose2020explaining, demetrio2019explaining}, and hierarchical
LIME~\cite{mitchell2024generating} explains opcode sequences. Aryal et al.~\cite{aryal2024explainability} sum MalConv's
SHAP values over the regions of each PE file and over 4\,KB subsections of its sections, use them to place evasive
perturbations, and find that averaging these region-level explanations over 500 files gives no usable global
explanation, so each file needs its own. FINER~\cite{he2023finer} formalizes the path from problem space to feature
space to explanation for security classifiers and sets EMBER-style detectors aside, arguing that their hashing and
histogram summaries lose the information an analyst would need from an explanation.

Where an application needs to cross the two spaces, the mapping is built by hand. Explanation-guided backdoors select
EMBER features by their SHAP values and realize them with a purpose-built PE editor, keeping only the 17 features that
could be mapped to file content~\cite{severi2021explanation}; explanation-guided evasion modifies features one at a time
with type-specific transformations~\cite{rosenberg2020generating}; and shortcut analyses explain detectors with a surrogate tree and review the top-ranked features by hand, finding
packing and metadata artifacts~\cite{mohammed2026beyond}. The
closest case in the other direction summarizes a raw-byte detector's saliency by overlaying it on PE structures parsed
from each file, which shows MalConv's reliance on the header~\cite{coull2019activation, bose2020explaining}; the
overlay is a hand-built, per-format map that sums each region's bytes, whereas our pushforward splits the credit of
shared bytes among the named features automatically, for any extractor. Detectors over different representations are
compared by their outputs~\cite{quiring2020against, salman2024robustness}, not by their evidence.

\subsection{Transferring Attributions}
\label{sec:XML:BgTransfer}

Redistributing an attribution from a stage's outputs to its inputs in proportion to each input's contribution is an
established idea. DeepLIFT's chain rule~\cite{shrikumar2017learning} and LRP~\cite{bach2015pixel} apply it layer by
layer inside a network, and G-DeepSHAP~\cite{chen2022explaining} lifts it to a series of models, propagating Shapley
values from one model to the next with a generalized rescale rule that satisfies efficiency at every stage. These
methods need every stage to be a model with its own attribution method; a feature extractor is a program, with
parsing, hashing, and windowed statistics over discrete, variable-length input and no natural baseline. In the other
direction, virtual inspection layers~\cite{vielhaben2024explainable} move relevance through an invertible linear
transform, such as time to frequency, into a domain the model never computes. Two other lines of work describe a model
in a vocabulary it does not compute: a global surrogate~\cite{craven1996extracting, frosst2017distilling} is trained to
imitate the model's outputs from interpretable inputs and is explained in its place, and concept-based
methods~\cite{bau2017network, kim2018interpretability} probe a network's internal activations for human-defined
concepts, each learned from a set of labeled examples. Grouping attributions over pixels into
segments or tokens into words~\cite{jullum2021groupshapley} is the simplest case, a 0/1 map within one space. Recording
which inputs produce an output is itself an old idea: provenance in databases~\cite{green2007provenance} and dynamic
taint tracking in programs do so; system provenance graphs for intrusion detection~\cite{milajerdi2019poirot,
cheng2024kairos} record causal relations among processes and files and are unrelated.

\subsection{Our Position}
\label{sec:XML:BgPosition}

Explanation provenance differs from this work in three ways. \emph{First}, it obtains the stage-wise matrix for a
feature extractor that is a program, by instrumenting it once, and so brings
attribution propagation to exactly the detectors FINER leaves out; the mapping that explanation-guided attacks and
shortcut analyses build by hand, feature by feature, becomes a by-product of extraction, for every feature and every
file. Our pullback coincides with G-DeepSHAP's rule, and we claim no novelty for that propagation step; what is new is
the matrix, which we ground in the Shapley value of a program's items and which instrumentation records exactly for
additive features. \emph{Second}, the same matrix runs the other way: the pushforward moves a byte-level
explanation through a non-invertible program into the features of a detector that never reads them, giving a raw-byte
detector a vocabulary in which its reliance recurs across files. We cast it as the inverse of the pullback and state
what any transfer can recover from a pulled-back explanation. Summing saliency over the parsed regions of a
file~\cite{coull2019activation, bose2020explaining, aryal2024explainability} moves explanations in the same direction,
but only into regions that do not share bytes; to our knowledge, no prior work transfers explanations into features
that share their elements while conserving attribution. Unlike a surrogate, the pushforward explains the detector itself, not a model that
imitates its outputs; unlike concept probes, it needs neither the detector's internals nor labeled examples, only a
byte-level explanation and the extractor's record of which bytes each feature is computed from. \emph{Third}, the two directions together put detectors over different representations
into one space, which turns explanation transfer into two security tasks: comparing two detectors on the same file, at a
resolution that windows of bytes do not reach, and auditing a raw-byte detector across a corpus in named features.
